% from ega1-9-fr.tex
\oldpage[I]{172}
comme homomorphisme de préfaisceaux~; tenant compte des définitions de
$\mathscr{G}$ et $\mathscr{H}$, il suffit donc, pour tout ensemble ouvert
$W\subset Y$, de définir un homomorphisme
\[
  u_W:\bigoplus_i\Gamma(f^{-1}(W)\cap X_i,\mathscr{F})
  \longrightarrow
  \bigoplus_{i,j}\Gamma(f^{-1}(W)\cap X_i\cap X_j,\mathscr{F})
\]
de façon à satisfaire aux conditions de compatibilité usuelles lorsque $W$
varie. Si, pour toute section
$s_i\in\Gamma(f^{-1}(W)\cap X_i,\mathscr{F})$, on désigne par $s_{i|j}$ sa
restriction à $f^{-1}(W)\cap X_i\cap X_j$, on posera
\[
  u_W((s_i))=(s_{i|j}-s_{j|i})
\]
et les conditions de compatibilité sont remplies de façon évidente. Pour
prouver que le noyau $\mathscr{R}$ de $u$ est $f_*(\mathscr{F})$, définissons
un homomorphisme de $f_*(\mathscr{F})$ dans $\mathscr{R}$ en faisant
correspondre à toute section $s\in\Gamma(f^{-1}(W),\mathscr{F})$ la famille
$(s_i)$, où $s_i$ est la restriction de $s$ à $f^{-1}(W)\cap X_i$~; les
axiomes (F 1) et (F 2) des faisceaux (G, II, 1.1) entraînent que cet
homomorphisme est \emph{bijectif}, ce qui termine la démonstration dans ce cas.

\item[b)] Dans le cas général, le même raisonnement s'applique une fois que
l'on a établi que les $\mathscr{F}'_{ij}$ sont quasi-cohérents. Or, par
hypothèse, $X_i\cap X_j$ est réunion finie d'ouverts affines $X_{ijk}$~; et
comme les $X_{ijk}$ sont des ouverts affines \emph{dans un schéma},
l'intersection de deux quelconques d'entre eux est encore un ouvert affine
(\hyperref[I.5.5.6-fr]{5.5.6}). On est donc ramené au premier cas, et
(9.2.1) est donc démontrée.
\end{enumerate}
\end{proof}

\begin{corollary}[9.2.2]
\label{I.9.2.2-fr}
La conclusion de (9.2.1) est valable dans chacun des cas suivants~:
\begin{enumerate}
\item[a)] $f$ est séparé et quasi-compact.
\item[b)] $f$ est séparé et de type fini.
\item[c)] $f$ est quasi-compact et l'espace sous-jacent à $X$ est localement
noethérien.
\end{enumerate}
\end{corollary}

\begin{proof}
Dans le cas a), les $X_{\alpha i}\cap X_{\alpha j}$ sont affines
(\hyperref[I.5.5.6-fr]{5.5.6}). Le cas b) est un cas particulier de a)
(\hyperref[I.6.6.3-fr]{6.6.3}). Enfin, dans le cas c), on peut se ramener au
cas où $Y$ est affine et l'espace sous-jacent à $X$ noethérien~; alors $X$
admet un recouvrement ouvert affine fini $(X_i)$, et les $X_i\cap X_j$, étant
quasi-compacts, sont réunions finies d'ouverts affines
(\hyperref[I.2.1.3-fr]{2.1.3}).
\end{proof}

\subsection{Prolongement des sections de faisceaux quasi-cohérents.}
\label{subsection:I.9.3-fr}

\begin{theorem}[9.3.1]
\label{I.9.3.1-fr}
Soit $X$ un préschéma dont l'espace sous-jacent est noethérien, ou un schéma
dont l'espace sous-jacent est quasi-compact. Soient $\mathscr{L}$ un
$\mathscr{O}_X$-Module inversible
(\hyperref[0.5.4.1-fr]{0, 5.4.1}), $f$ une section de $\mathscr{L}$ au-dessus
de $X$, $X_f$ l'ensemble ouvert des $x\in X$ où $f(x)\neq0$
(\hyperref[0.5.5.1-fr]{0, 5.5.1}), $\mathscr{F}$ un
$\mathscr{O}_X$-Module quasi-cohérent.
\begin{enumerate}
\item[(i)] Si $s\in\Gamma(X,\mathscr{F})$ est telle que $s|X_f=0$, il existe
un entier $n>0$ tel que $s\otimes f^{\otimes n}=0$.
\item[(ii)] Pour toute section $s\in\Gamma(X_f,\mathscr{F})$, il existe un
entier $n>0$ tel que $s\otimes f^{\otimes n}$ se prolonge en une section de
$\mathscr{F}\otimes\mathscr{L}^{\otimes n}$ au-dessus de $X$.
\end{enumerate}
\end{theorem}

\begin{proof}
\begin{enumerate}
\item[(i)] Comme l'espace sous-jacent à $X$ est quasi-compact, donc réunion
finie d'ouverts affines $U_i$ tels que $\mathscr{L}|U_i$ soit isomorphe à
$\mathscr{O}_X|U_i$, on est ramené au cas où $X$ est affine et
$\mathscr{L}=\mathscr{O}_X$. Dans ce cas, $f$ s'identifie à un élément de
$A(X)$ et on a $X_f=D(f)$~; $s$ s'identifie à un élément d'un
$A(X)$-module $M$ et $s|X_f$ à l'élément correspondant de $M_f$, et le
résultat est trivial, compte tenu de la définition d'un module de fractions.
